% Generated by D:/tools/scratch_qdd/mmrec/render.py from data_d.py - edit the data, not this file.
\subsection{D. VLA --- 세부 조종 스틱이 결정을 따르게, 그리고 끝-끝으로}\label{rec:D}
\noindent\textbf{한 줄:} 움직임 줄만 채택(+0.032)되고 MolmoAct식 보조·카메라 3대·층별 KV·명령 입력은 모두 불채택. 결정 토큰을 청크가 따르지 않는 문제(E-SR0 0.34--0.43)는 시뮬에서 먼 구간 반사실 분기로 0.991까지 풀었지만 실데이터는 0.536(E-SR1e 진행 중). 확신도는 기록만. 다음은 ser-A-min-3 재학습과 끝-끝 학습(E-VLA-E2E).

\paragraph{D1. 배경: VLM 사용법, 소규모 끝-끝 준비, 하나로 융합}
{\small 09-24 23:05 -- 09-25 23:05 KST}
\begin{itemize}
  \item \textbf{사용자} (ul 54, 09-24 23:05 KST): \intent{어떤식으로 vlm을 써야할지를 좀 생각해봐야할듯 엔드투 엔드 목적으로 사용하게 될거니까 하긋ㅂ할때 어케하게될지}
  \item \textbf{사용자} (ul 54, 09-24 23:05 KST): \intent{A$\rightarrow$B 단계적 (권고)}
  \item \textbf{그래서(결정):} 정본 §51: 먼저 결정층(보기 토큰)만 학습, 뒤에 같은 백본에 행동 전문가를 붙여 전체 끝-끝으로. §55 단계 A SFT: A\_SFT $-$ A\_zero +0.504 [+0.498, +0.511], p95 0.307 s.
  \item \textbf{사용자} (ul 56, 09-25 02:13 KST): \intent{사실 이건 실험인거고 제대로 된 데이터 를 제대로 넣은 실험은 아닌거 아는거지? 지금한건 소규모 테스트로 실제 동작하는지를 본거지?}
  \item \textbf{그래서(결정):} §56: §55는 파이프라인 확인용 소규모 시험 --- 단계 3 수치 대부분에 같은 성격 표시.
  \item \textbf{사용자} (ul 59--61, 02:55--02:59 KST): \intent{그냥 계속 다 자동으로 해주면돼 아직 엔드투엔드는 아니고 모든 파이프라인 전체를 엔드투엔드 학습할 준비까지 모두 준비를 마쳐주면 되는거야 알겠지? 검증해보면서 계속 자동으로}
  \item \textbf{사용자} (ul 59--61, 02:55--02:59 KST): \intent{각 모듈별로 따로 놀면서ㅗ 가각의 성능에 의존하는 파이프라인보다는 최대한 하나로 융합되는게 좋은 연구인거알지?}
  \item \textbf{사용자} (ul 59--61, 02:55--02:59 KST): \intent{어쩔수없는건 물론 그냥 하긴해야하구}
  \item \textbf{사용자} (ul 59--61, 02:55--02:59 KST): \intent{이렇게 모듈 각각 검증해보고 다하고나면 한번 돌아다니면서 객관 검증해보고 2연속 사이클 도는데 문제가 없으면, 이제 다음인 소규모 엔드투 엔드 준비해서 데이터 오픈소스 된것중에 해서 한 10gb정도로만해서 소규모 엔드투엔드 학습한번해보자.}
  \item \textbf{그래서(결정):} §58: 런타임은 하나의 융합 모델(Qwen3-VL-4B 결정 + 흐름 행동 전문가). R7 객관 검증 순회 $\rightarrow$ 2회 연속 무결이 E2E 준비 기준점.
  \item \textbf{사용자} (ul 66, 04:18 KST): \intent{직므은 단순이니까 10으로 하되 나중에는 우리가 데이터하고 할떄는 30으로 모을거}
  \item \textbf{결과:} S-E2E(ROBOTIS 공개 실물 RB1 + RB2, 원본 8.95 GB, 10 Hz $\rightarrow$ 42,813행, train 37,484 / val 1,799): 학습 PASS(bc1bd54, 20:58, dec 4.623 $\rightarrow$ 0.681·0.688), 진단 D1--D3(a4a6257: 계산 한계 아님, 라벨은 학습 가능·일반화 간극, 모호성 아님), 규모 곡선(3478cf2, 23:05) 1k/9.4k/18.7k/37.5k = 0.552/0.682/0.686/0.683 --- 같은 스텝에서 데이터 레버 약하고 포화.
  \item \textbf{그래서(결정):} 다음 레버는 입력(시간 맥락) 쪽(D2).
\end{itemize}

\paragraph{D2. 시간 맥락 $\rightarrow$ 움직임 줄 채택, 2프레임 비디오 불채택}
{\small 09-25 21:42 -- 09-26 02:24 KST}
\begin{itemize}
  \item \textbf{해 봄:} 조사 118def0(21:42): 과거 프레임 1장을 Qwen3-VL 2프레임 비디오로(1순위), 움직임 줄 + 드롭아웃.
  \item \textbf{사용자} (ul 71, 21:48--21:49 KST): \intent{1순위 아이디가 좋다잉}
  \item \textbf{사용자} (ul 71, 21:48--21:49 KST): \intent{저거는 지금 당장 오버리프에 업뎃해줘도 좋겠는데}
  \item \textbf{사용자} (ul 71, 21:48--21:49 KST): \intent{피규어나 이런거에도}
  \item \textbf{해 봄:} E-TC 2$\times$2 사전 등록(13:05:07Z, 0b853c2 22:07): 2프레임 비디오 V $\times$ 움직임 줄 M, se2e, 2,000스텝, 검증 300.
  \item \textbf{결과:} f4f6a49(09-26 00:39): 네 칸 0.683 / 0.696 / 0.708 / 0.721; M 주효과 +0.025 [+0.003, +0.047], 전환 층 +0.027, FULL p95 $-$3.8 \%; V +0.013(< 0.02), FULL p95 +20.4 \%(CPU 전처리).
  \item \textbf{사용자} (ul 80, 00:40 KST): \intent{일단 움직임 줄 채택할게}
  \item \textbf{그래서(결정):} 정본 §83(f00cb59, 00:40): 움직임 줄 \texttt{se2e-\allowbreak{}motion@v1} 채택, 2프레임 비디오 불채택. 발견된 누수: S-E2E 속도가 중앙 차분(미래 0.1 s 포함) $\rightarrow$ 인과 후방 차분 \texttt{se2e\_\allowbreak{}c1}(44c907d, 00:53).
  \item \textbf{결과:} 확인 실험(7dced6a, 02:24): 고친 데이터·시드 2개·검증 1,799에서 합동 +0.032 [+0.022, +0.043], 전이 +0.030 $\rightarrow$ CONFIRMED, §83 유지.
\end{itemize}

\paragraph{D3. MolmoAct 1차 이식 --- 모두 불채택}
{\small 09-26 01:41 -- 10:22 KST}
\begin{itemize}
  \item \textbf{사용자} (ul 83, 01:41--01:42 KST): \intent{VLA도 사진을 보는것도 좋을 것 같구}
  \item \textbf{사용자} (ul 83, 01:41--01:42 KST): \intent{엔드포인트의 움직임을 위에 덧그리는 방식도 뭔가 도움이 될 수 있을 것 같기도 하구}
  \item \textbf{결과:} E-MA1(06355f0 등록) 관문 G0 실패(ee1b000, 02:41): 머리 카메라 FK 투영 명목 100.9/264.5 px, PnP 뒤 86.9/212.1 px(문턱 12/30) $\rightarrow$ 학습 0칸, R2 뒤 재등록.
  \item \textbf{결과:} E-MA1b 로봇 기준 3D 미래 손끝 궤적 보조(ded38d7 등록): 검증 1,799 합동 +0.0031 [$-$0.0022, +0.0085], 전이 0.000 $\rightarrow$ 불채택(2df9bdc, 04:49; §85 보충).
  \item \textbf{결과:} E-CAM3 머리 + 양 손목(시각 토큰 356 $\rightarrow$ 460, 5c46056): +0.0014 [$-$0.0041, +0.0069], FULL 결정 p95 +11.7 \%(84.6 $\rightarrow$ 94.6 ms) $\rightarrow$ 불채택(76ab10c, 07:22); 결정층 입력은 머리 + 활성 손목 유지, 양 손목은 Astra 요청에만(§57 보충).
  \item \textbf{결과:} E-MA3 층별 KV 조건(40895f8): 청크 오차 합동 상대 감소 +2.19 \% [+1.48 \%, +2.88 \%] < 문턱 5 \% $\rightarrow$ 불채택(76ab10c; §84 보충 3).
  \item \textbf{결과:} E-MA2 VLA 입력에 Astra \texttt{edit}(fde4f6a 등록, R2): 돌린 명령 준수율 문장 C1 0.177·화살표 C2 0.001(둘 다 < 0.6); C1은 결정층이 명령 칸을 0.97로 고르지만 청크는 0.28(90$^\circ$)·0.09(180$^\circ$)만 따름(399d1cb, 10:22).
  \item \textbf{그래서(결정):} §84 보충 4: 융합 VLA 입력에 Astra 명령을 넣지 않고 설계 §11 코드 편향만. 결정 토큰이 청크를 조종한다고 가정하지 않는다($\rightarrow$ D5).
\end{itemize}

\paragraph{D4. MolmoAct가 `교훈'으로만 작동 $\rightarrow$ 준비부터, 그리고 실데이터로만}
{\small 09-26 11:25 -- 19:38 KST}
\begin{itemize}
  \item \textbf{사용자} (ul 95, 11:25--11:26 KST): \intent{나는 몰모엑트가 우리가 제대로 준비를 안해놔서 도움이 안되는 교훈ㅇ,로 작동한듯}
  \item \textbf{사용자} (ul 95, 11:25--11:26 KST): \intent{준비를 아ㅖ제대로해서 검증 필요. 시간 걸려두되니까}
  \item \textbf{결과:} R2 준비 관문(6cae25c, 12:04, 학습 없음): 성립 조건 14개; G-cam 표지 중심 0.21 px, G-fk URDF FK 대 시뮬 손끝 9.5 mm 실패, G-trace 현재점 영상 안 0.994, 만들 것 M1--M9.
  \item \textbf{사용자} (ul 99, 14:29 KST): \intent{시뮬거쳤다오는건 절대 반대 몰모엑트 참고해}
  \item \textbf{사용자} (ul 99, 14:29 KST): \intent{추가학습·MolmoAct 이식만 실데이터 (Recommended)}
  \item \textbf{그래서(결정):} 정본 §89(b8c231f, 14:29): 추가 학습·MolmoAct 이식은 실데이터 + 실영상 모델 라벨로, R2 시뮬 MolmoAct 데이터 생성은 중단(파일럿 8편 뒤, 기본 경로 비트 동일 5/5 확인); 조건 준수 연구·폐루프 시험장으로서 시뮬은 유지.
  \item \textbf{결과:} 실데이터 준비(0ca87c5, 15:26): Molmo2-ER 두 팔 포인팅 10,834회, 원시 오류 27--30 \%(대부분 다른 팔); 사전 고정 거르개 v1 실패(유지 0.18) $\rightarrow$ v2 남긴 점 오류 0.037·유지 0.945; 같은 장면·다른 경로 시연은 없음 $\rightarrow$ 조종 데이터(S)는 실물 수집 필요.
  \item \textbf{해 봄:} E-MAR-real 요인 A(7245787, 15:46; 변경 1 3e2f5f6): RB2 실데이터에 MolmoAct 2D 궤적 보조 \texttt{trace5-\allowbreak{}point@v1}.
  \item \textbf{결과:} 825fd37(19:38): 검증 1,799 합동 +0.0047 [$-$0.0006, +0.0103], 전이 +0.0026, RB2 라벨 행 $-$0.003; 보조 궤적 오차 62 px(평균 궤적 기준선 145 px) --- 궤적은 배웠지만 결정 정확도로 옮지 않음, 준수 불변(0.341 대 0.340). 약 3.6 GPU-h.
  \item \textbf{그래서(결정):} §84 보충 11: 불채택. 실데이터 MolmoAct 이식은 대체 경로 시연을 모은 뒤 E-MAR-S로만, 이력 덧그림 O는 보류.
\end{itemize}

\paragraph{D5. 조이스틱 준수 문제 발견 $\rightarrow$ 목표 0.8}
{\small 09-26 10:37 -- 11:22 KST}
\begin{itemize}
  \item \textbf{해 봄:} E-SR0(7b42b97 등록, 학습 없음): 기존 체크포인트에서 전문가에 주는 확정 결정만 반사실로 바꿔 청크 FK 변위가 따르는지(cos > 0.5).
  \item \textbf{결과:} §84 보충 5(10:52 KST): S-E2E 0.340(우연 0.335)·R2 0.431(우연 0.333) $\rightarrow$ WEAK; 참 결정일 때 R2 청크 0.92, 이산 \texttt{phase\_\allowbreak{}id}는 그리퍼 0.95로 따름.
  \item \textbf{사용자} (ul 91, 10:59 KST): \intent{x2점유시작}
  \item \textbf{사용자} (ul 91, 10:59 KST): \intent{준수율이 0.8은 될방법론알 찾아야해}
  \item \textbf{그래서(결정):} 결합은 결정 토큰이 아니라 청크·손끝 기준에 직접 적용(§84 보충 5). E-SR1b 채택 문턱 0.6 $\rightarrow$ 0.8. 방법론 조사(f6c7c3f, 11:22): 1순위 먼 구간 반사실 분기 데이터, 2 권한 가중 구조, 3 CFG w(a), 4 권한 게이트 adaLN; 거리 권한 a(C3)로 모두 재평가.
\end{itemize}

\paragraph{D6. E-SR1b 드롭아웃 + CFG·되짚기 $\rightarrow$ 도달 못 함}
{\small 09-26 11:31 -- 13:51 KST}
\begin{itemize}
  \item \textbf{해 봄:} 사전 등록 ad85725(11:31): R2, E-MA2 C0 레시피에 결정 드롭아웃 0.3 + 결정 CFG(w 1--8)·되짚기 결정 라벨·둘 다, 시드 2개씩 8판.
  \item \textbf{결과:} 12f5902(13:51): 반사실 준수 A\_xy 최고 0.598(되짚기 B, 기준 0.438)·그럴듯한 편집 P\_xy 최고 0.398(기준 0.358); B는 청크 MSE +14.6 \%·그리퍼 일치 $-$0.030으로 비열등 실패; 드롭아웃은 w = 1 준수를 0.362로 낮춤; CFG w 5--8은 청크 MSE +106--357 \%. 약 7.9 GPU-h.
  \item \textbf{그래서(결정):} §84 보충 7 NOT\_REACHED: 드롭아웃 + CFG는 준수 수단 후보에서 뺀다, 결합은 청크·손끝 기준 경유 유지.
\end{itemize}

\paragraph{D7. E-SR1c 먼 구간 반사실 분기 $\rightarrow$ 채택(시뮬)}
{\small 09-26 12:09 -- 14:29 KST}
\begin{itemize}
  \item \textbf{해 봄:} 사전 등록 a79bd64(12:09; 변경 1 2a2ad34, 변경 2 beacc0d): R2 먼 구간(대상까지 $\geq$ 10 cm) 스냅숏에 강제 조이스틱 결정 4개와 그 방향 IK 직선 청크를 짝지은 반사실 분기 128,661행(Isaac 재생 관문 G-br: 추종 중앙 0.9 mm·방향 100 \%·침투 0)을 전문가 배치의 50 \%로; C2는 권한 게이트 adaLN 추가.
  \item \textbf{결과:} e091312(14:29): 먼 구간 반사실 방향 준수 C1 0.991(C0 0.344; A\_z 0.989·$\rho$ 0.992, 분기 없는 bottle\_tray 0.979, 편집 $\pm$30·45$^\circ$ 0.976), 근접 비열등 N1--N5·전 구간 (i)--(iv) 통과, 근접 반사실도 0.868로 따름. 대가: 먼 구간 참 결정 청크 MSE +29 \%·끝점 +4 mm. C2 adaLN 이득 없음($-$0.006)·지연 +9 ms. 약 2.5 GPU-h.
  \item \textbf{그래서(결정):} §84 보충 8 ADOPT\_C1: 먼 구간(a = 1)만 조이스틱 결정 경유 조종 허용, 근접(a = 0)은 VLA 자신의 결정만; §6 투영·§11 편향 $\times$ a; 융합 VLA 레시피에 먼 구간 분기 50 \%. 구현은 Task 17(C5).
\end{itemize}

\paragraph{D8. E-SR1d 같은 레시피를 실데이터로 $\rightarrow$ 도달 못 함}
{\small 09-26 14:59 -- 16:09 KST}
\begin{itemize}
  \item \textbf{해 봄:} 사전 등록 32065de(14:59): S-E2E \texttt{se2e\_\allowbreak{}c1}에 시뮬 없이 URDF IK 기구학 분기 61,993행(관문 G-kin: 거부 11.3 \%, 한계·속도 위반 0), 권한 a = 그리퍼 사건 되짚기 대용(관문 G-a 통과, 참 거리 검증 불가).
  \item \textbf{결과:} 2c93c9b(16:09): 먼 구간 A\_xy 0.339 $\rightarrow$ 0.536(+0.198 [+0.187, +0.209], 시드 0.552·0.521), A\_z 0.519 $\rightarrow$ 0.772, 근접 비열등 통과. 약 2.2 GPU-h.
  \item \textbf{그래서(결정):} §84 보충 10 NOT\_REACHED: 실데이터 체크포인트에는 먼 구간 결정 경유 조종을 옮기지 않는다(청크·손끝 경유 유지); §92 P1(실데이터 준수 $\geq$ 0.8)은 불성립.
\end{itemize}

\paragraph{D9. E-SR1e --- 원인은 적합 부족, 두 팔로 다시}
{\small 09-26 18:59 KST \textasciitilde{} (진행 중)}
\begin{itemize}
  \item \textbf{결과:} 학습 없는 진단(등록 0.1절): E-SR1d C1의 TRAIN 먼 구간 준수 0.549·0.523 = 검증 0.552·0.521 $\rightarrow$ 일반화 차가 아니라 적합 부족; 모든 검증 과제가 학습에 있음, 편별 고르게 `반만' 따름; 2,000스텝 = 0.43 에폭에서 손실이 아직 내려감; 실데이터 전문가는 결정 토큰을 거의 안 씀(C0 참 결정 적중 0.47).
  \item \textbf{해 봄:} 사전 등록 7b076cd(18:59): 팔 A(6,000스텝·분기 50 \%)·D(3,000스텝·분기 75 \%) + L1 대조 결정 안내(E-SR1d C1, w 1.5·2·3, 먼 구간만); 코드 기준 = E-SR1d 등록 커밋(HEAD의 6질문 판은 옛 체크포인트를 못 읽음); RB3는 팔에 넣지 않음(일반화 차 없음).
  \item \textbf{진행 중·대기:} \tentative{결과 문서(\texttt{results/\allowbreak{}sr1e.\allowbreak{}md})는 아직 없다 --- handoff §0의 `진행 중' 줄(메인 GPU 2·3 학습 $\rightarrow$ 평가). 결과에 따라 실데이터 레시피·P1이 바뀐다(ADOPT이면 §84 보충 10 수정).}
\end{itemize}

\paragraph{D10. SigLIP --- 필요 없으면 하지 않는다}
{\small 09-26 10:43 -- 11:03 KST}
\begin{itemize}
  \item \textbf{사용자} (ul 90, 10:43--10:49 KST): \intent{Siglip을 쓸 방법도 찾아봐 어떻게 엮을 수 있을지}
  \item \textbf{사용자} (ul 90, 10:43--10:49 KST): \intent{나는 저걸 우리vla에 쓰는게 어떨가 싶은데}
  \item \textbf{사용자} (ul 90, 10:43--10:49 KST): \intent{그래c,e반영가능성도 검사해보고 결정하자 일단ema2 다되고 나면 그거보고 ㄱ}
  \item \textbf{사용자} (ul 92, 11:03 KST): \intent{만약 시그립이 그래서 필요없다면 굳이 더하란 내말은 무시해두돼}
  \item \textbf{결과:} 조사(4e2e242): 융합 VLA의 영상 탑(Qwen3-VL-4B)이 이미 SigLIP2-Large 계열(24층·폭 1024·패치 16, 동결).
  \item \textbf{그래서(결정):} §84 보충 6(14d78b5, 11:03): 별도 SigLIP 통합은 하지 않음; 손목 패치 (e0)·두 번째 SigLIP 줄 (c)는 준수 문제를 푼 뒤 이득 근거가 있을 때만(기본 보류).
\end{itemize}

\paragraph{D11. 새로움 점수와 결정 확신도 $\rightarrow$ 게이트 없이 기록만}
{\small 09-26 14:36 -- 20:07 KST}
\begin{itemize}
  \item \textbf{해 봄:} E-NOV0(d225a84, 14:36): 백본 특징 k-NN 새로움 점수(학습 없음)로 모델 자신의 큰 오류·새 환경을 미리 아는가.
  \item \textbf{결과:} 7c26faf(15:27): 큰 오류 AUROC 0.570, 90 \% 재현율 호출 감소 13 \%, dr 이동 0.742·과제 빼기 0.751 $\rightarrow$ NONE(흐름 게이트·T\_nov 채널 안 둠, §84 보충 9). 판정 밖: 결정 확신도 0.804(감소 0.450).
  \item \textbf{사용자} (ul 106, 15:37 KST): \intent{결정 확신도는 어쩌다가 발견했음?}
  \item \textbf{사용자} (ul 106, 15:37 KST): \intent{확신도 기반으로 해두면 좋을듯 왜 그런 선택을 했는지도.}
  \item \textbf{사용자} (ul 106, 15:37 KST): \intent{그리고 조금씩 확신도를 능력을 더 높이는 개선도 필요할듯}
  \item \textbf{그래서(결정):} 정본 §95(780ca10, 15:40): 확신도(질문별 1·2등 로그확률 차의 최솟값) 기반 운용 + 선택 이유 기록 + 개선 트랙 --- 결과 뒤 선택이므로 새 표본에서 먼저 검증(E-CONF).
  \item \textbf{해 봄:} E-CONF(6191369, 16:18; 변경 1 f37f8f8·393a184).
  \item \textbf{결과:} cd29570(19:13): R2-새(학습된 적 없는 R2\_TRAIN 실패 편 244편·8,208) AUROC 0.806 [0.796, 0.817]·감소 0.379 통과, S-E2E 실데이터 0.748 [0.724, 0.771]·0.244 실패(문턱 0.75·0.30); 개선 후보 \texttt{temp\_\allowbreak{}joint} +0.041/+0.024(문턱 +0.03 미달).
  \item \textbf{그래서(결정):} §95 보충 1 NONE: 흐름 우선 표시·보수 모드의 트리거로 쓰지 않고 결정마다 `선택 이유' 기록만(훅 \texttt{conf-\allowbreak{}base@v1}, 5a940d2 20:07, 어떤 제어 경로도 읽지 않음).
\end{itemize}

\paragraph{D12. Astra를 추가 학습 교사로(§88) $\rightarrow$ 실데이터 E-AT 설계}
{\small 09-26 14:27 -- 14:46 KST}
\begin{itemize}
  \item \textbf{사용자} (ul 98, 14:27 KST): \intent{아스트라도 추가학습의 일환으로 써야겠음}
  \item \textbf{사용자} (ul 98, 14:27 KST): \intent{본학습은 가정한 지티로 하돼 추가학습으로 실제 반영하는거지}
  \item \textbf{그래서(결정):} 정본 §88(bc923c2, 14:27): 본 학습 = 가정한 정답(규칙 라벨·시뮬 오라클), 추가 학습 = 결합 폐루프에서 Astra가 실제로 낸 판단·수정(DAgger식). 선행 조건: Astra 명령 정확도 측정(E-ACC) --- 틀린 라벨을 배우지 않게.
  \item \textbf{결과:} E-AT 설계(cb81017, 14:46, 문서만): ul 99에 따라 실데이터 전용 --- RB2 실물 시연에 Astra 라벨(구간 품질·J2 진단 위주), 비용 계획 20,000원·상한 25,000원.
  \item \textbf{진행 중·대기:} \tentative{E-AT는 실행 없음: 유료라 사용자 승인 필요(ul 114), \texttt{quality:} 줄은 \texttt{ser-\allowbreak{}A-\allowbreak{}min-\allowbreak{}3}에 없음(결합 계획 G8 --- \texttt{ser-\allowbreak{}A-\allowbreak{}min-\allowbreak{}4}에 묶기).}
\end{itemize}

\paragraph{D13. 직렬화 ser-A-min-3 판 올림 $\rightarrow$ 재학습 대기}
{\small 09-26 15:17 -- 09-27 00:45 KST}
\begin{itemize}
  \item \textbf{해 봄:} 결합 Task 12 한 번의 판 올림: 움직임 줄 + 그리퍼 질문(340c5dc, 15:17), \texttt{prompt\_\allowbreak{}config} 판본·범주 기록과 검사(a761b2b), PH-A1 문구(b931f75), 구간이 그리퍼 사건을 품음·학습 때만 unknown 구간 p 0.3·프롬프트 원천 해시 보강(bce0bd8, 18:45), 재개 시험(fd24d2e, 00:45).
  \item \textbf{결과:} 새 판은 기존 체크포인트를 모두 거부한다 $\rightarrow$ 폐루프·파드 검증이 모두 막힘(C8).
  \item \textbf{진행 중·대기:} \tentative{[결정 필요] ser-A-min-3 재학습(결합 계획 D2 추천: VLA-sim 먼저, VLA-real은 E-SR1e 뒤). book 01 VLA-R3는 [예정].}
\end{itemize}

\paragraph{D14. 우리 VLA는 흔한 VLA가 아니다, 그래도 끝-끝으로 학습}
{\small 09-27 01:4x KST}
\begin{itemize}
  \item \textbf{사용자} (ul 128, 01:4x KST): \intent{나는 VLA 부분은 엔드투 엔드로 학습이 될 필요가 있다고 생각해. 하지만 알겠지만 저 VLA가 우리가 아는 VLA의 역할이 아니라는건 인지하고 있지?}
  \item \textbf{그래서(결정):} 정본 §96 보충 4(a54170f, 01:39): 역할 = 상위 의도(구간 줄) + 영상·상태를 받아 typed 결정(방향·크기·대상·단계 + 쥐기·놓기 시점)을 내는 빠른 조이스틱, 과제를 스스로 해석하지 않음. E2E(E-VLA-E2E, 무료, 새 등록): 결정은 감독된 중간 출력으로 남기고 KI식 전문가 $\rightarrow$ 백본 절연 유지 + 이산 행동 토큰(E1), 정답 $\rightarrow$ 예측 결정 조건 교과(E2), 구간 줄 의도 교란(E4); 결합 고리 안의 VLA DAgger(본학습 자리만).
  \item \textbf{진행 중·대기:} \tentative{선행: ser-A-min-3 단계식 재학습(비교 기준)과 E-SR1e 결과. 행동 토큰은 \texttt{ser-\allowbreak{}A-\allowbreak{}min-\allowbreak{}4} 판 올림이라 E-AT \texttt{quality:} 줄과 한 번에 묶는다. [예정].}
\end{itemize}
